\chapter{Monopoles, Dyons, and Quantum Sectors}
\label{ch:monopoles-dyons-quantum-sectors}
\ConventionsLorentzian

The smooth monopole core of
Chapter~\ref{ch:vol2-classical-yang-mills-matter} is only the local part of
the physics.  A magnetic particle is also a long-range charged sector.  Its
asymptotic electromagnetic field changes the rotation generators, the
allowed two-particle wavefunctions, and the meaning of electric charge in a
\(\theta\)-vacuum.  This chapter isolates those sector-level facts before
any attempt is made to use a monopole moduli-space metric as a quantum
Hamiltonian.

The guiding rule is that a monopole calculation has two inputs.  The core
and its collective coordinates determine a short-distance completion and the
low-velocity coordinates.  The fields at infinity determine the charge
lattice, the angular momentum stored outside the cores, and the Hilbert space
on which scattering or bound-state questions are posed.  A physical monopole
chapter has to keep both inputs visible.

\section{Dyonic Charges and the DSZ Pairing}
\label{sec:dyonic-charges-dsz-pairing}

Work in a low-energy Abelian frame with primitive magnetic charge
\(\mathfrak g_{\rm m}\) and primitive electric charge \(e_\infty\), normalized
by
\[
  e_\infty\mathfrak g_{\rm m}=2\pi .
\]
A dyonic sector is labelled by
\(\gamma=(n_{\mathrm e},n_{\mathrm m})\in\mathbb Z^2\).  In a
\(\theta\)-vacuum the electric field measured at infinity is shifted by the
Witten effect:
\begin{equation}
  Q_{\mathrm E}(\gamma;\theta)
  =
  e_\infty
  \left(
    n_{\mathrm e}
    +
    \frac{\theta n_{\mathrm m}}{2\pi}
  \right),
  \qquad
  Q_{\mathrm M}(\gamma)
  =
  \mathfrak g_{\rm m}n_{\mathrm m}.
  \label{eq:monopole-chapter-witten-shifted-charges}
\end{equation}
The important invariant for two charged sectors is not either electric
charge separately.  It is the antisymmetric pairing
\begin{equation}
  \langle\gamma_1,\gamma_2\rangle_{\rm DSZ}
  :=
  \frac{
    Q_{\mathrm E}(\gamma_1;\theta)Q_{\mathrm M}(\gamma_2)
    -
    Q_{\mathrm M}(\gamma_1)Q_{\mathrm E}(\gamma_2;\theta)
  }{2\pi}.
  \label{eq:monopole-chapter-dsz-physical-pairing}
\end{equation}

\begin{proposition}[DSZ pairing]
\label{prop:monopole-chapter-dsz-pairing}
For \(\gamma_i=(n_{\mathrm e,i},n_{\mathrm m,i})\),
\begin{equation}
  \langle\gamma_1,\gamma_2\rangle_{\rm DSZ}
  =
  n_{\mathrm e,1}n_{\mathrm m,2}
  -
  n_{\mathrm m,1}n_{\mathrm e,2}
  \in\mathbb Z .
  \label{eq:monopole-chapter-dsz-integer-pairing}
\end{equation}
In particular, the pairwise locality condition is independent of
\(\theta\).
\end{proposition}

\begin{proof}
Substitute \eqref{eq:monopole-chapter-witten-shifted-charges} into
\eqref{eq:monopole-chapter-dsz-physical-pairing}.  The factor
\(e_\infty\mathfrak g_{\rm m}=2\pi\) cancels the denominator.  The two
\(\theta\)-terms are
\[
  \frac{\theta}{2\pi}n_{\mathrm m,1}n_{\mathrm m,2}
  -
  \frac{\theta}{2\pi}n_{\mathrm m,1}n_{\mathrm m,2},
\]
so they cancel before any quantization condition is imposed.  What remains is
the integral symplectic pairing of the charge labels.
\end{proof}

This cancellation is a diagnostic.  A formula for a monopole sector which
assigns physical importance to the separate real number
\(Q_{\mathrm E}Q_{\mathrm M}/2\pi\) has not yet formed the two-sector
locality invariant.  The \(\theta\)-dependent electric field is real, but the
Dirac-Schwinger-Zwanziger obstruction to putting two sectors in the same
quantum theory is the antisymmetric pairing
\eqref{eq:monopole-chapter-dsz-integer-pairing}.

\section{Field Angular Momentum and Monopole Harmonics}
\label{sec:field-angular-momentum-monopole-harmonics}

Two separated dyons carry angular momentum in their long-range fields.  For a
relative separation vector \(\mathbf R\), the electromagnetic field
contribution is
\begin{equation}
  \mathbf J_{\rm field}
  =
  \frac{
    Q_{\mathrm E}(\gamma_1;\theta)Q_{\mathrm M}(\gamma_2)
    -
    Q_{\mathrm M}(\gamma_1)Q_{\mathrm E}(\gamma_2;\theta)
  }{4\pi}\,
  \widehat{\mathbf R}
  =
  \frac12
  \langle\gamma_1,\gamma_2\rangle_{\rm DSZ}\,
  \widehat{\mathbf R}.
  \label{eq:monopole-chapter-field-angular-momentum}
\end{equation}
Thus a charge-monopole pair with odd DSZ pairing carries a half-integer unit
of angular momentum before any spin carried by the cores is included.

\begin{proposition}[Monopole harmonics]
\label{prop:monopole-chapter-monopole-harmonics}
Let \(N=\langle\gamma_1,\gamma_2\rangle_{\rm DSZ}\).  The angular
wavefunctions of the two-body relative coordinate are monopole harmonics with
\begin{equation}
  j=\frac{|N|}{2}+\ell,
  \qquad
  \ell=0,1,2,\ldots ,
  \label{eq:monopole-chapter-allowed-angular-momenta}
\end{equation}
and \(m=-j,-j+1,\ldots,j\).  The radial Hamiltonian contains the angular
barrier
\begin{equation}
  \frac{1}{2\mu r^2}
  \left[
    j(j+1)-\frac{N^2}{4}
  \right],
  \label{eq:monopole-chapter-monopole-harmonic-barrier}
\end{equation}
where \(\mu\) is the reduced mass.  For \(N=0\) this reduces to the ordinary
spherical-harmonic barrier \(\ell(\ell+1)/(2\mu r^2)\).
\end{proposition}

\begin{proof}
The vector potential of one dyon is a connection on the line bundle over the
linking two-sphere seen by the other dyon.  Its first Chern number is the DSZ
integer \(N\).  Sections of this line bundle replace ordinary functions on
\(S^2\).  The covariant rotation algebra has irreducible representations
starting at \(j=|N|/2\), and tensoring with ordinary orbital excitations gives
\eqref{eq:monopole-chapter-allowed-angular-momenta}.  The covariant angular
operator is the total angular momentum with the radial field contribution
removed, giving the eigenvalue
\(j(j+1)-N^2/4\) in the relative radial problem.
\end{proof}

This is already a quantum statement about monopoles before one solves for any
smooth core.  The core removes the Dirac singularity and supplies the
short-distance completion.  The long-range DSZ integer supplies the angular
momentum bundle of the two-particle problem.  Treating monopoles as ordinary
scalar particles with ordinary \(Y_{\ell m}\) partial waves would erase the
field angular momentum in \eqref{eq:monopole-chapter-field-angular-momentum}.

\section{BPS Alignment as a No-Force Benchmark}
\label{sec:bps-alignment-no-force-benchmark}

In a BPS embedding, the Abelian fields at infinity are accompanied by a
massless scalar tail.  The no-force property of mutually aligned BPS dyons is
not a statement about the dimension of a moduli space.  It is a cancellation
between vector repulsion and scalar attraction in the long-distance
Hamiltonian.

Use the two-component charge vector
\[
  \mathbf q_i=(Q_{\mathrm E,i},Q_{\mathrm M,i})
\]
and let the scalar charge of a positive-energy BPS particle on the same
central-charge ray be \(|\mathbf q_i|\).  In this Abelian BPS-tail
normalization the leading static interaction is
\begin{equation}
  V_{\rm tail}(R)
  =
  \frac{
    \mathbf q_1\cdot\mathbf q_2
    -
    |\mathbf q_1|\,|\mathbf q_2|
  }{4\pi R}.
  \label{eq:monopole-chapter-bps-tail-potential}
\end{equation}

\begin{controlledapproximation}[Aligned BPS no force]
\label{ca:monopole-chapter-aligned-bps-no-force}
If \(\mathbf q_2=\lambda\mathbf q_1\) with \(\lambda>0\), then
\[
  V_{\rm tail}(R)=0.
\]
If the charge vectors are not on the same positive ray, the tail
\eqref{eq:monopole-chapter-bps-tail-potential} is nonzero unless additional
light fields or channel projections are included.

For \(\mathbf q_2=\lambda\mathbf q_1\), the dot product is
\(\lambda|\mathbf q_1|^2\), while
\(|\mathbf q_1|\,|\mathbf q_2|=\lambda|\mathbf q_1|^2\).  The two terms in
\eqref{eq:monopole-chapter-bps-tail-potential} cancel.  Conversely, equality
in the Cauchy--Schwarz inequality with positive sign is exactly the condition
that the two nonzero charge vectors lie on the same positive ray.
\end{controlledapproximation}

This benchmark is useful because it says what a monopole moduli-space
approximation is allowed to mean physically.  A flat center-of-mass metric is
not enough.  The asymptotic fields must also make the long-range force vanish
in the BPS channel, and the two-body Hilbert space must still carry the DSZ
angular momentum data of Section~\ref{sec:field-angular-momentum-monopole-harmonics}.
When the BPS alignment or supersymmetric protection is absent, the same
classical magnetic core can receive long-range forces, radiation, and
nonzero-mode corrections which have to be included before a particle
spectrum, scattering phase, or bound-state claim is made.
